\documentclass[../../../include/open-logic-section]{subfiles}

\begin{document}

\olfileid{his}{set}{mythology}

\olsection{迷思多于史实？}

带着一个多世纪的后见之明回顾这些事件，我们必须谨慎，不可盲从这些评判。
这些结果无疑是新颖、令人振奋且出人意料的。但它们到底有多么令人震惊呢？
它们是否真的证明了我们不应依赖几何直观？

至于所谓的震惊，\citet{Gouvea2011}指出 康托尔给 戴德金
的那句著名便条“\emph{je le vois, mais je ne le crois
pas}”（我看到了，但我不相信）多少有些断章取义。
以下是更多的上下文（引自 \citeauthor{Gouvea2011}）：

\begin{quote}
请原谅我对此主题的热忱，倘若我如此麻烦您的善意与耐心；
我最近寄给您的那些发现，即使对我自己来说也如此出乎意料、如此新奇，
以至于在我从您——敬爱的朋友——那里获得关于其正确性的判定之前，
我无法心安。只要您还未认可我的结论，我就只能说：
\emph{je le vois, mais je ne le crois pas.}（我看到了，但我不相信。）
\end{quote}

康托尔 知道他的结果“如此出乎意料、如此新奇”。
不过，他是否真的觉得自己的结果\emph{不可信}，是很值得怀疑的。
正如 \citeauthor{Gouvea2011} 所指出的，他仅仅是请 戴德金 核验他所提供的证明。

至于几何直观：皮亚诺发表他的空间填充曲线时并未附上任何图表。
但当 希尔伯特发表他的曲线时，他解释了自己的意图：
只要读者们“借助如下的几何直观”，他就能为他们提供一种理解 皮亚诺 结果的清晰途径；
随后他附上了一系列\emph{图表}，就如同第 \olref[his][set][pathology]{sec} 节中所提供的那样。

更一般地说：尽管图表在已发表的证明中已不再流行，但无法回避的事实是，
数学家们在进行证明时\emph{经常}使用图表。（粗略地说：优秀的数学家知道何时可以依赖几何直观。）

简而言之：不要轻信那些夸大之词；或者至少，不要不加甄别地盲信。关于此话题的更多内容，你可以阅读 \citet{Giaquinto2007}。

\end{document}